Chương này trình bày các cơ sở lý thuyết và kiến thức nền tảng cần thiết cho đề tài. Nội dung bao gồm: tổng quan về phát hiện gian lận thẻ tín dụng, giới thiệu về thuật toán XGBoost, các kỹ thuật feature engineering trong bài toán fraud detection, và cuối cùng là các khái niệm về drift trong machine learning.

\section{Tổng quan về phát hiện gian lận thẻ tín dụng}

\subsection{Bài toán phát hiện gian lận}

Phát hiện gian lận thẻ tín dụng (\textit{Credit Card Fraud Detection}) là bài toán phân loại nhị phân trong đó mỗi giao dịch được phân loại là một trong hai nhãn: hợp pháp (legitimate) hoặc gian lận (fraud). Đây là một bài toán đặc biệt trong machine learning vì nó có những đặc điểm riêng biệt:

\begin{itemize}
    \item \textbf{Mất cân bằng dữ liệu nghiêm trọng}: Trong hầu hết các bộ dữ liệu fraud detection, tỷ lệ gian lận chỉ chiếm khoảng 0.1-0.5\% tổng số giao dịch. Điều này khiến các mô hình phân loại truyền thống có xu hướng bỏ qua lớp thiểu số (fraud) để đạt accuracy cao.
    
    \item \textbf{Hành vi gian lận thay đổi theo thời gian}: Kẻ gian lận liên tục thay đổi chiến lược để tránh bị phát hiện, dẫn đến phân phối của các mẫu gian lận thay đổi theo thời gian.
    
    \item \textbf{Chi phí sai lệch không đồng nhất}: Chi phí của việc bỏ sót gian lận (false negative) cao hơn nhiều so với việc đánh dấu nhầm giao dịch hợp pháp là gian lận (false positive).
    
    \item \textbf{Yêu cầu về thời gian thực}: Hệ thống cần phát hiện gian lận trong thời gian thực hoặc gần thực để ngăn chặn giao dịch gian lận.
\end{itemize}

Các phương pháp tiếp cận bài toán phát hiện gian lận có thể được chia thành các nhóm chính sau:

\begin{enumerate}
    \item \textbf{Phương pháp dựa trên quy tắc (Rule-based)}: Sử dụng các quy tắc được định nghĩa thủ công để phát hiện gian lận. Ví dụ: nếu số tiền giao dịch lớn hơn ngưỡng nào đó và giao dịch diễn ra ở quốc gia khác với địa chỉ thẻ thì đánh dấu là nghi ngờ.
    
    \item \textbf{Phương pháp học có giám sát (Supervised Learning)}: Sử dụng các thuật toán machine learning với dữ liệu đã được gắn nhãn để huấn luyện mô hình phân loại. Các thuật toán phổ biến bao gồm: Logistic Regression, Random Forest, XGBoost, \ac{lstm}.
    
    \item \textbf{Phương pháp học không giám sát (Unsupervised Learning)}: Phát hiện các giao dịch bất thường mà không cần dữ liệu gắn nhãn. Ví dụ: Isolation Forest, One-Class SVM, Autoencoder.
    
    \item \textbf{Phương pháp kết hợp (Hybrid)}: Kết hợp nhiều phương pháp để tận dụng ưu điểm của từng phương pháp.
\end{enumerate}

Trong đề tài này, chúng tôi sử dụng phương pháp học có giám sát với thuật toán XGBoost, kết hợp với các kỹ thuật xử lý mất cân bằng dữ liệu.

\subsection{Bộ dữ liệu Credit Card Fraud}

Bộ dữ liệu được sử dụng trong đề tài là Credit Card Fraud Detection Dataset từ Kaggle, được thu thập và công bố trong nghiên cứu của \cite{creditcard_dataset}. Bộ dữ liệu này chứa các giao dịch thẻ tín dụng của người dùng châu Âu trong vòng 2 ngày vào tháng 9 năm 2013.

Bộ dữ liệu có các đặc điểm sau:
\begin{itemize}
    \item \textbf{Time}: Thời gian tính bằng giây kể từ giao dịch đầu tiên trong bộ dữ liệu.
    \item \textbf{V1-V28}: 28 đặc trưng được tạo thông qua \ac{pca} để bảo mật thông tin khách hàng.
    \item \textbf{Amount}: Số tiền giao dịch.
    \item \textbf{Class}: Nhãn phân loại (1 = fraud, 0 = legitimate).
\end{itemize}

Tổng số giao dịch: 284,807, trong đó có 492 giao dịch gian lận (0.173\%).

\section{Giới thiệu về thuật toán XGBoost}

\subsection{Tổng quan về Gradient Boosting}

XGBoost (eXtreme Gradient Boosting) là một thuật toán thuộc họ Gradient Boosting, được phát triển bởi Chen và Guestrin vào năm 2016 \cite{xgboost}. Thuật toán này xây dựng mô hình theo phương pháp ensemble, kết hợp nhiều cây quyết định yếu (\textit{weak learners}) để tạo ra một mô hình mạnh.

Gradient Boosting hoạt động bằng cách xây dựng các cây quyết định theo chuỗi, trong đó mỗi cây mới được huấn luyện để sửa lỗi của các cây trước đó. Cụ thể, tại mỗi bước, thuật toán tính gradient của hàm mất mát đối với dự đoán hiện tại, sau đó huấn luyện một cây mới để fit gradient này.

Hàm mất mát được sử dụng trong XGBoost cho bài toán phân loại nhị phân là log loss:

\begin{equation}
L(\phi) = \sum_{i=1}^{n} y_i \log(p_i) + (1-y_i) \log(1-p_i)
\end{equation}

trong đó $y_i$ là nhãn thực và $p_i$ là xác suất dự đoán.

\subsection{Các đặc điểm nổi bật của XGBoost}

XGBoost có một số đặc điểm nổi bật khiến nó trở thành lựa chọn phổ biến trong các bài toán phân loại:

\begin{enumerate}
    \item \textbf{Hiệu suất cao}: XGBoost sử dụng các kỹ thuật tối ưu hóa như parallel processing, cache-aware access, và compressed column blocks để tăng tốc độ huấn luyện.
    
    \item \textbf{Xử lý overfitting}: XGBoost có các tham số điều chỉnh như \textit{max\_depth}, \textit{min\_child\_weight}, \textit{subsample}, \textit{colsample\_bytree} để kiểm soát overfitting.
    
    \item \textbf{Xử lý missing values**: XGBoost có khả năng tự động xử lý các giá trị missing mà không cần xử lý trước.
    
    \item \textbf{Hỗ trợ cross-validation**: XGBoost tích hợp sẵn các phương pháp cross-validation để đánh giá mô hình.
    
    \item \textbf{Khả năng mở rộng**: XGBoost có thể chạy trên nhiều nền tảng và hỗ trợ các ngôn ngữ lập trình khác nhau (Python, R, Java, C++).
\end{enumerate}

Trong bài toán phát hiện gian lận, XGBoost có những ưu điểm sau:
- Xử lý tốt dữ liệu mất cân bằng thông qua tham số \textit{scale\_pos\_weight}.
- Khả năng feature importance cao, giúp hiểu được các đặc trưng quan trọng nhất.
- Tốc độ huấn luyện nhanh, phù hợp với các bộ dữ liệu lớn.

Các tham số quan trọng của XGBoost trong bài toán fraud detection:

\begin{itemize}
    \item \textit{n\_estimators}: Số lượng cây trong ensemble.
    \item \textit{max\_depth}: Độ sâu tối đa của mỗi cây.
    \item \textit{learning\_rate}: Tốc độ học (eta).
    \item \textit{scale\_pos\_weight}: Trọng số cho lớp positive (fraud).
    \item \textit{early\_stopping\_rounds}: Số vòng không cải thiện trước khi dừng sớm.
    \item \textit{eval\_metric}: Metric đánh giá trong quá trình huấn luyện.
\end{itemize}

\section{Kỹ thuật Feature Engineering}

Feature engineering đóng vai trò quan trọng trong việc cải thiện hiệu suất của mô hình machine learning. Trong bài toán phát hiện gian lận, các đặc trưng được tạo ra dựa trên hiểu biết về bản chất của giao dịch gian lận.

\subsection{Đặc trưng về thời gian}

Các giao dịch gian lận thường có mối liên hệ với thời gian. Chúng tôi tạo ra các đặc trưng thời gian sau:

\begin{itemize}
    \item \textbf{time\_diff}: Khoảng thời gian giữa giao dịch hiện tại và giao dịch trước đó. Giá trị này được giới hạn trong khoảng [0, 86400] (tối đa 1 ngày).
    
    \item \textbf{time\_diff\_log}: Log của time\_diff để giảm ảnh hưởng của các giá trị ngoại lai.
\end{itemize}

Công thức tính:
\begin{equation}
\text{time\_diff}_t = \text{Time}_t - \text{Time}_{t-1}
\end{equation}
\begin{equation}
\text{time\_diff\_log}_t = \log(1 + \text{time\_diff}_t)
\end{equation}

Lý do: Các gian lận thường xảy ra với tần suất cao trong thời gian ngắn (ví dụ: nhiều giao dịch nhỏ liên tiếp để test thẻ), do đó time\_diff nhỏ có thể là dấu hiệu của gian lận.

\subsection{Đặc trưng về số tiền}

Số tiền giao dịch là một trong những yếu tố quan trọng trong phát hiện gian lận. Chúng tôi tạo ra các đặc trưng sau:

\begin{itemize}
    \item \textbf{log\_amount}: Log của số tiền giao dịch để giảm ảnh hưởng của các giao dịch có số tiền lớn.
    
    \item \textbf{amt\_to\_mean\_ratio}: Tỷ lệ số tiền giao dịch so với số tiền trung bình trong dữ liệu huấn luyện.
    
    \item \textbf{amt\_to\_median\_ratio}: Tỷ lệ số tiền giao dịch so với số tiền trung vị.
    
    \item \textbf{is\_high\_amount}: Nhị phân, đánh dấu các giao dịch có số tiền lớn hơn ngưỡng phân vị 95\%.
\end{itemize}

Công thức tính:
\begin{equation}
\text{log\_amount} = \log(1 + \text{Amount})
\end{equation}
\begin{equation}
\text{amt\_to\_mean\_ratio} = \frac{\text{Amount}}{\mu_{\text{Amount}}}
\end{equation}
\begin{equation}
\text{amt\_to\_median\_ratio} = \frac{\text{Amount}}{\text{median}_{\text{Amount}}}
\end{equation}
\begin{equation}
\text{is\_high\_amount} = \mathbb{1}(\text{Amount} > P_{95})
\end{equation}

Lưu ý: Khi tính các đặc trưng này cho dữ liệu test hoặc production, cần sử dụng các thống kê từ dữ liệu huấn luyện để tránh data leakage.

\section{Khái niệm về Drift trong Machine Learning}

Drift là hiện tượng phân phối dữ liệu hoặc mối quan hệ giữa đặc trưng và nhãn thay đổi theo thời gian. Trong các hệ thống machine learning thực tế, drift là một vấn đề phổ biến và ảnh hưởng nghiêm trọng đến hiệu suất của mô hình.

\subsection{Các loại Drift}

Có ba loại drift chính:

\begin{enumerate}
    \item \textbf{Data Drift (Covariate Shift)}: Là sự thay đổi trong phân phối của các đặc trưng đầu vào ($P(X)$), trong khi phân phối của nhãn ($P(Y|X)$) không thay đổi.
    
    \textit{Ví dụ}: Nếu ban đầu khách hàng chủ yếu ở độ tuổi 25-35, sau đó khách hàng chuyển sang độ tuổi 35-45, thì đây là data drift.
    
    \item \textbf{Concept Drift}: Là sự thay đổi trong mối quan hệ giữa đặc trưng và nhãn ($P(Y|X)$). Nói cách khác, quy luật để phân loại nhãn đã thay đổi.
    
    \textit{Ví dụ}: Trong fraud detection, hành vi gian lận thay đổi theo thời gian, các giao dịch trước đây được coi là bình thường có thể trở thành gian lận và ngược lại.
    
    \item \textbf{Prediction Drift}: Là sự thay đổi trong phân phối của các dự đoán của mô hình. Đây có thể là dấu hiệu của cả data drift và concept drift.
\end{enumerate}

Hình \ref{fig:drift_types} minh họa ba loại drift:

\begin{figure}[H]
    \centering
    \includegraphics[width=0.8\textwidth]{img/drift_types.png}
    \caption{Các loại Drift trong Machine Learning}
    \label{fig:drift_types}
\end{figure}

\subsection{Phương pháp phát hiện Drift}

Có nhiều phương pháp để phát hiện drift, được chia thành hai nhóm chính:

\begin{enumerate}
    \item \textbf{Phương pháp thống kê (Statistical Methods)}:
    \begin{itemize}
        \item \textbf{Kolmogorov-Smirnov (KS) Test}: So sánh phân phối của hai tập dữ liệu.
        \item \textbf{Population Stability Index (PSI)}: Đo lường sự khác biệt giữa hai phân phối.
        \item \textbf{Chi-squared test}: Kiểm tra xem hai tập dữ liệu có cùng phân phối hay không.
    \end{itemize}
    
    \item \textbf{Phương pháp dựa trên error rate}:
    \begin{itemize}
        \item Theo dõi precision, recall, F1 trên dữ liệu mới có nhãn.
        \item So sánh performance giữa các thời điểm khác nhau.
    \end{itemize}
\end{enumerate}

Trong đề tài này, chúng tôi sử dụng KS test và PSI để phát hiện data drift, và theo dõi F1-score để phát hiện concept drift.

PSI là một chỉ số phổ biến trong lĩnh vực banking và credit scoring. Công thức tính PSI:

\begin{equation}
\text{PSI} = \sum_{i=1}^{n} (p_i - q_i) \log\left(\frac{p_i}{q_i}\right)
\end{equation}

trong đó $p_i$ là tỷ lệ phần trăm của dữ liệu hiện tại trong bucket $i$, và $q_i$ là tỷ lệ phần trăm của dữ liệu tham chiếu trong bucket $i$.

Ngưỡng PSI thường được sử dụng:
\begin{itemize}
    \item PSI < 0.1: Không có drift đáng kể.
    \item 0.1 $\le$ PSI < 0.25: Drift nhẹ, cần theo dõi.
    \item PSI $\ge$ 0.25: Drift nghiêm trọng, cần huấn lại mô hình.
\end{itemize}

\subsection{Xử lý Drift}

Khi phát hiện drift, có một số chiến lược xử lý:

\begin{enumerate}
    \item \textbf{Huấn lại định kỳ (Periodic Retraining)}: Huấn lại mô hình theo lịch cố định (ví dụ: hàng tuần, hàng tháng).
    
    \item \textbf{Huấn lại khi phát hiện drift (Triggered Retraining)}: Huấn lại mô hình khi drift được phát hiện vượt ngưỡng.
    
    \item \textbf{Online Learning}: Cập nhật mô hình liên tục khi có dữ liệu mới.
    
    \item \textbf{Ensemble Methods}: Kết hợp nhiều mô hình, sử dụng mô hình phù hợp nhất cho từng thời điểm.
\end{enumerate}

Trong đề tài này, chúng tôi sử dụng chiến lược triggered retraining kết hợp với kiểm tra định kỳ.

\section{Tổng kết chương}

Chương 2 đã trình bày các cơ sở lý thuyết cần thiết cho đề tài, bao gồm:

\begin{itemize}
    \item Tổng quan về bài toán phát hiện gian lận thẻ tín dụng và đặc điểm của bộ dữ liệu.
    \item Giới thiệu về thuật toán XGBoost và các đặc điểm nổi bật của nó.
    \ Các kỹ thuật feature engineering được sử dụng trong đề tài.
    \item Khái niệm về drift và các phương pháp phát hiện drift.
\end{itemize}

Các kiến thức này sẽ được áp dụng trong Chương 3 để xây dựng hệ thống hoàn chỉnh.